\chapter{Normalization Methods}

\section{Batch Normalization}

Training deep neural networks can be sensitive to the scale and distribution of activations. Normalization methods address this issue by standardizing intermediate representations.

\medskip

\noindent
\textbf{Batch normalization (BatchNorm).}  
Given activations $z$ in a layer, BatchNorm transforms them as follows:
\[
\hat{z}_i = \frac{z_i - \mu_B}{\sqrt{\sigma_B^2 + \varepsilon}},
\]
where $\mu_B$ and $\sigma_B^2$ are the mean and variance computed over a mini-batch.

\medskip

\noindent
\textbf{Affine transformation.}  
After normalization, a learnable transformation is applied:
\[
y_i = \gamma \hat{z}_i + \beta,
\]
where $\gamma$ and $\beta$ are parameters.

\medskip

\noindent
\textbf{Purpose.}
\begin{itemize}
    \item Stabilize the distribution of activations
    \item Improve gradient flow
    \item Accelerate training
\end{itemize}

\medskip

\noindent
\textbf{Interpretation.}  
Batch normalization reduces sensitivity to the scale of parameters and inputs.

\medskip

\noindent
\textbf{Training vs inference.}
\begin{itemize}
    \item During training: statistics computed from mini-batches
    \item During inference: running averages are used
\end{itemize}

\medskip

\noindent
\textbf{Practical effect.}
\begin{itemize}
    \item Allows larger learning rates
    \item Reduces dependence on initialization
\end{itemize}

\section{Layer and Instance Normalization}

While BatchNorm relies on batch statistics, alternative normalization methods operate independently of the batch.

\subsection{Layer Normalization}

\medskip

\noindent
\textbf{Definition.}  
Layer normalization normalizes across features within a single data point:
\[
\hat{z} = \frac{z - \mu_{\text{layer}}}{\sqrt{\sigma_{\text{layer}}^2 + \varepsilon}}.
\]

\medskip

\noindent
\textbf{Properties.}
\begin{itemize}
    \item Independent of batch size
    \item Suitable for sequence models (e.g., transformers)
\end{itemize}

\subsection{Instance Normalization}

\medskip

\noindent
\textbf{Definition.}  
Instance normalization normalizes each sample independently, often across spatial dimensions in vision tasks.

\medskip

\noindent
\textbf{Applications.}
\begin{itemize}
    \item Image generation
    \item Style transfer
\end{itemize}

\medskip

\noindent
\textbf{Comparison of methods.}

\begin{itemize}
    \item BatchNorm: uses batch statistics
    \item LayerNorm: uses feature statistics per sample
    \item InstanceNorm: uses spatial statistics per sample
\end{itemize}

\medskip

\noindent
\textbf{Insight.}  
Different normalization methods reflect different assumptions about the structure of the data.

\section{Stabilization of Training}

Normalization methods improve the stability and efficiency of training.

\medskip

\noindent
\textbf{Effect on optimization.}
\begin{itemize}
    \item Smooths the loss landscape
    \item Reduces sensitivity to parameter scaling
    \item Improves conditioning
\end{itemize}

\medskip

\noindent
\textbf{Gradient flow.}
\begin{itemize}
    \item Prevents vanishing and exploding gradients
    \item Maintains consistent signal propagation
\end{itemize}

\medskip

\noindent
\textbf{Internal covariate shift (informal).}  
As parameters change during training, the distribution of activations shifts. Normalization reduces this effect.

\medskip

\noindent
\textbf{Regularization effect.}
\begin{itemize}
    \item BatchNorm introduces noise through batch statistics
    \item Can improve generalization
\end{itemize}

\medskip

\noindent
\textbf{Learning rate interaction.}
\begin{itemize}
    \item Enables use of larger learning rates
    \item Leads to faster convergence
\end{itemize}

\medskip

\noindent
\textbf{Insight.}  
Normalization reshapes the optimization problem, making it easier to solve.

\section{Theoretical and Practical Insights}

Although normalization methods are widely used, their theoretical understanding is still evolving.

\medskip

\noindent
\textbf{Scale invariance.}  
Normalization makes the network less sensitive to scaling of parameters.

\medskip

\noindent
\textbf{Reparameterization.}  
Normalization can be viewed as changing the parameterization of the model, which affects optimization dynamics.

\medskip

\noindent
\textbf{Geometry of optimization.}  
Normalization effectively modifies the geometry of the loss landscape, often making it more well-conditioned.

\medskip

\noindent
\textbf{Empirical observations.}
\begin{itemize}
    \item Improves training speed
    \item Enhances stability across architectures
    \item Essential in many modern deep networks
\end{itemize}

\medskip

\noindent
\textbf{Limitations.}
\begin{itemize}
    \item BatchNorm depends on batch size
    \item May behave poorly with very small batches
\end{itemize}

\medskip

\noindent
\textbf{Design considerations.}
\begin{itemize}
    \item Choice of normalization depends on architecture and data
    \item Often combined with other techniques (e.g., residual connections)
\end{itemize}

\medskip

\noindent
\textbf{Insight.}  
Normalization is both a practical tool and a conceptual mechanism for improving optimization.

\section{A Unifying Perspective}

Normalization methods can be understood as transformations that:

\begin{itemize}
    \item standardize representations,
    \item improve conditioning,
    \item stabilize gradient flow,
    \item introduce implicit regularization.
\end{itemize}

\medskip

\noindent
\textbf{Core components.}
\begin{itemize}
    \item Normalization step
    \item Learnable scaling and shifting
    \item Interaction with optimization dynamics
\end{itemize}

\medskip

\noindent
\textbf{Key insight.}  
Normalization plays a central role in making deep networks trainable at scale.

\section{Outlook}

This chapter examined normalization methods in deep learning:
\begin{itemize}
    \item batch normalization and its variants,
    \item alternative normalization strategies,
    \item their role in stabilizing training,
    \item theoretical and practical perspectives.
\end{itemize}

\medskip

In the next chapter, we study regularization techniques specific to deep learning, including dropout and data augmentation.

\medskip

\noindent
\textbf{Guiding principle.}  
Well-designed transformations of intermediate representations can significantly improve both optimization and generalization.